\section{Persistent \texorpdfstring{$E_8$}{E8} Information}

\paragraph{Modeling status.} The following identifications of lattice labels with spin, color and charge are proposed physical correspondences. A representation-theoretic construction connecting them to Lorentz and gauge transformations is still required.

Before discussing particles, their decay rates, or their lifetimes, we must define what persists. The $E_8$ substrate is an active, discrete network. To exist over time, structured information must successfully transit from one lattice node to an adjacent node during each fundamental clock cycle. 

To survive this transition, the data payload must satisfy the Universal Architecture of Persistence: it must respect the node's finite channel budget (\textbf{Capacity}), preserve its causal history (\textbf{Map}), synchronize phase frames with the destination (\textbf{Protocol}), remain within local hardware bounds (\textbf{Governor}), pay the irreducible thermodynamic cost of transition (\textbf{Toll}), and maintain its signal above the resolution floor (\textbf{Margin}).

If a state violates any of these pillars during transit, it triggers Causal Aliasing, the collapse of distinct causal histories into an indistinguishable superposition, resulting in immediate dissolution into the vacuum noise floor. 

The familiar properties of elementary particles emerge strictly at the intersections of these architectural requirements:

\begin{itemize}
    \item \textbf{Map + Protocol $\rightarrow$ Phase Offset (Spin):} The Map requires a reconstructible causal history, and the Protocol requires lossless phase synchronization across discrete links. Their intersection forces a strict geometric phase-tracking requirement.
    \item \textbf{Protocol + Governor $\rightarrow$ Internal Embedding (Color):} The Protocol mediates interactions, but the Governor bounds the local structural complexity. Their intersection forces the isolation of internal coordinate spaces to prevent capacity overflow.
    \item \textbf{Map + Capacity $\rightarrow$ Boundary Flux (Charge):} The Map defines the local identity of the state, and the Capacity defines the macroscopic causal grid. Their intersection dictates exactly how much of the internal payload must anchor to the external manifold.
    \item \textbf{Toll and Margin:} These operate at two scales. At the substrate level, they are fixed geometric invariants: the minimum cost of any state transition, and the minimum signal distinguishable from noise. For specific routing protocols, these scale to protocol-specific latency and solvency thresholds. A protocol's lifetime is determined by the competition between its accumulated maintenance Toll and the Transition Toll required to reach a lower-latency ground state. We derive this explicitly in Section \cref{sec:particle_decay}.
\end{itemize}

A ``particle,'' therefore, is an information packet executing a specific routing protocol that successfully maintains all six pillars during node-to-node transfer. 

\subsection{Causal Synchronization and the Phase Offset (spin)}

When an information packet moves across a discrete lattice link, the \textit{Protocol} pillar (Paper I) mandates that it must maintain phase coherence with the destination node to preserve probability (Unitarity). 

To achieve this, the transition operator must align the payload's local coordinate frame. As established in the substrate derivation, the $\nu=16$ active capacity channels are encoded on the chiral state spaces of the $D_4$ sub-lattices. Because these are chiral spaces, their geometry strictly operates on the double-cover of the orthogonal rotation group ($Spin(8)$). 

This hardware architecture forces an irreducible geometric mismatch during data transfer. Traversing a closed spatial loop on the single-cover coordinate manifold translates to a standard $2\pi$ phase advance. However, returning the chiral payload (on the double-cover) to its exact identity state requires a full $4\pi$ rotation. Consequently, for every step the spatial frame takes, the internal payload only completes half of its phase cycle. 

We define this irreducible geometric phase lag as the \textbf{Causal Synchronization Offset}. 

In standard physics, this exact $1:2$ geometric ratio translates to \textbf{Spin-1/2}. It is the necessary phase-tracking metadata required to synchronize a chiral payload across a discrete link without triggering Landauer erasure (loss of Unitarity).

This synchronization mechanism imposes a strict functional limit on the network. By Rule 1: Capacity Partitioning (Paper III), a node possesses a hard limit of exactly $\nu=16$ active channels. If two information packets attempt to occupy the exact same channel while carrying the exact same Synchronization Offset, the network's transition operator cannot distinguish between their updates. 

Because the node lacks the residual bandwidth to resolve this ambiguity, the overlapping payloads trigger \textbf{Causal Aliasing} (Paper II). The local capacity overflows, causing the states to instantly decohere and dissolve into the vacuum noise floor. 

This strict algorithmic requirement for collision avoidance on a finite-capacity network physically manifests as the \textbf{Pauli Exclusion Principle}. Fermionic statistics are thus derived not as an axiomatic quantum rule, but as the mandatory routing logic required to prevent catastrophic data collisions on the $E_8$ substrate.


\subsection{The Interaction Remainder and Internal Links (color \& gluons)}

By the rule of Capacity Partitioning (established in \textit{Nested Persistence}\cite{meyer_ienested_2026}), the $E_8$ substrate allocates a finite bandwidth to node-to-node interactions, defined by the interaction rank $\sigma = 5$. However, to manifest as a persistent, localized entity (a distinct Map), an information packet must anchor itself using the topological boundary of a 2-sphere ($\chi = 2$). 

This geometric requirement strictly partitions the interaction capacity. Subtracting the topological boundary from the total interaction rank yields an isolated 3-dimensional subspace, which we define as the \textbf{Interaction Remainder}:
\begin{equation}
    N_{res} = \sigma - \chi = 5 - 2 = 3
\end{equation}
As established in Paper I\cite{meyer_informational_2026}, these 3 degrees of freedom represent the internal spatial embedding—the coordinate address space required to orient the 2D boundary within the local substrate. 

Information must flow across this internal address space to maintain structural coherence. According to the \textit{Entropic Action} rules (Paper II), this requires a 1-form connective link (a network Protocol) that synchronizes these 3 internal embedding coordinates between adjacent nodes. Crucially, this specific link operates \textit{strictly} within the $N_{res} = 3$ subspace, mediating the exact capacity saturation threshold we previously derived as the Strong Coupling ($\alpha_s$) in Paper III.

Mapping this to standard physics the 3-dimensional Interaction Remainder is the \textbf{Color Bulk}, and the 1-form transition link that synchronizes its embedding coordinates is the \textbf{Gluon}. 

This purely architectural derivation structurally dictates the confinement of the strong force. In the PM framework, the macroscopic $D=4$ spacetime manifold only registers information that projects \textit{across} the $\chi=2$ topological boundary. An internal link (the gluon) merely permutes the 3D embedding address of that boundary; it does not alter the boundary itself. Because it strictly routes data behind the topological interface, it generates no external Map update. The macroscopic manifold cannot ``see'' it. Therefore, the gluon is structurally isolated from the macroscopic observer, yielding exactly the phenomenon of color confinement without requiring an independent physical force mechanism.

\subsection{Map and Capacity: Boundary Flux and the Topological Footprint (charge)}

If the internal embedding links (the Color Bulk) operate entirely orthogonally to the macroscopic manifold, how does the internal state interact with the external universe? 

For a localized data state to possess a persistent identity (a \textbf{Map}) while propagating across the macroscopic spatial grid (\textbf{Capacity}), it must establish a geometric anchor between its internal payload and the external $D=4$ spacetime. As established in Paper I, the only valid geometric anchor on this substrate is the topological boundary of a 2-sphere ($\chi=2$). 

As the packet moves from node to node, a specific fraction of its total information payload must be processed \textit{through} this $\chi=2$ topological interface to maintain its macroscopic anchor. We define the \textbf{Boundary Flux} as this exact geometric footprint: the ratio of the payload's capacity that permeates the spacetime boundary versus the capacity retained in the internal interaction bulk. 

We identify this Boundary Flux physically as \textbf{Electric Charge ($Q$)}. 

In this framework, electric charge is not an axiomatic coupling parameter inserted into a Lagrangian; it is a strict measure of geometric permeability. It quantifies the exact topological footprint required for a Map to maintain its Capacity allocation while transiting across the continuous spacetime projection.